\subsection{Related work}
\label{sec:related-work}

We survey prior work on proving statements involving ring arithmetic, structured linear algebra, and lattice-based relations, organized by the approach taken to handle the ring structure.

%%%%%%%%%%%%%%%%%%%%%%%%%%%%%%%%%%%%%%%%%%%%%%%%%%%%%%%%%%%%%%%%%%%%%%%%%%%%%%%
\paragraph{Unrolling ring arithmetic into field circuits.}
The most direct approach is to decompose each ring element $a \in \ring_q$ into its $n$ coefficients over $\field_p$ and encode ring multiplication as an explicit NTT butterfly circuit.  Boschini, Camenisch, Ovsiankin, and Spooner~\cite{PQCRYPTO:BCOS20} instantiate this with the Aurora IOP~\cite{EC:BCRSVW19}, proving RSIS and RLWE relations by compiling them into R1CS over $\field_p$; the same strategy underlies vCNN~\cite{EPRINT:LKKO20} and ZEN~\cite{EPRINT:FQZDC21}. The fundamental limitation is that each ring multiplication contributes $O(n \log n)$ constraints, so an inner product $\langle a, s \rangle = t$ with $M$ ring elements yields $O(Mn \log n)$ constraints regardless of the backend proof system used.

%%%%%%%%%%%%%%%%%%%%%%%%%%%%%%%%%%%%%%%%%%%%%%%%%%%%%%%%%%%%%%%%%%%%%%%%%%%%%%%
\paragraph{Sumcheck-based FFT and convolution proofs.}
Liu, Xie, and Zhang~\cite{CCS:LiuXieZha21} observe that the DFT entry $\omega^{jm}$ has a multilinear extension that factors as a product of $\log n$ terms over the binary digits of $j$ and $m$. This lets a single sumcheck verify a length-$n$ FFT with $O(n)$ prover work, faster than computing the transform, and the verifier evaluates the factored MLE in $O(\log n)$. A ring convolution is proved by composing FFT, Hadamard-product, and inverse-FFT sumchecks inside a GKR framework~\cite{JACM:GolKalRot15}.  These can be amortized: the $M$ forward transforms are batched into one GKR sumcheck, and likewise for the $N$ inverse transforms. The dominant prover cost $O(NMn)$ matches ours; the verifier cost $O(\log(NMn))$ is lower, since each batched instance costs $O(\log n)$.

Our protocol folds the CRT matrix $\tilde{C}$ directly into the virtual polynomial of a single sumcheck expressing $As = t$, rather than composing separate transform sub-protocols via GKR.  Because $\tilde{C}$ enters as a factor, the verifier encounters it exactly once, evaluating $\tilde{C}(\alpha', \gamma)$ in $O(n)$; this covers all $NM$ ring products and does not scale with the matrix dimensions, giving verifier cost $O(n + \log(NM))$.

Before arriving at this solution, we developed a coefficient-domain sumcheck that proves ring multiplication directly as a negacyclic circulant matrix-vector product, bypassing the CRT decomposition entirely and recovering $O(\log n)$.  However, in the coefficient domain the convolution entangles the column-index and coefficient-index variables: the first sumcheck must range over all $\mu + \nu$ of them simultaneously, producing folding challenges that fix the oracle evaluation points at distinct last-$\nu$ coordinates. In~\cref{sec:speedyspartan}, we sketch a SpeedySpartan for ring arithmetic and show how shared evaluation coordinates yield a consolidation of these claims essentially for free. Concretely, these PCS optimizations save the verifier significantly more time than a single $O(n)$ evaluation takes, considering $n$ generally ranges from $64$ to $512$ for protocols at the $128$-bit security level.

%%%%%%%%%%%%%%%%%%%%%%%%%%%%%%%%%%%%%%%%%%%%%%%%%%%%%%%%%%%%%%%%%%%%%%%%%%%%%%%
\paragraph{SNARKs for ring arithmetic.}
Jia, Li, Xing, Yao, and Yuan~\cite{C:JLXYY25} establish proximity gaps for Reed--Solomon codes over Galois rings $\mathrm{GR}(2^k, d)$ and build a FRI-based polynomial commitment scheme, yielding the first publicly verifiable SNARKs for $\ZZ_{2^k}$-arithmetic with polylogarithmic proof size.  Their machinery targets power-of-two moduli, not the NTT-friendly setting we consider. Similarly, Wei, Zhang, and Deng~\cite{PKC:WeiZhaDen25} construct transparent SNARKs over arbitrary Galois rings by extending Brakedown's~\cite{C:GLSTW23} linear-time encodable codes to the ring setting, achieving $O(n)$ prover time, $O(\sqrt{n})$ proof size, and $O(\sqrt{n})$ verifier time for log-space uniform circuits.

These Galois-ring works address a complementary problem to ours.
They build general-purpose proof systems whose native computation domain is a ring $\ZZ_{2^k}$ (or its Galois extension), motivated by matching machine-word arithmetic or MPC-over-$\ZZ_{2^k}$ protocols.
Our work targets a different setting: proving structured linear-algebraic relations (matrix-vector products) over \emph{cyclotomic} rings $\ZZ_p[X]/(X^n+1)$ with splitting primes, as arise in lattice-based cryptography.
We exploit the CRT factorization specific to this setting to obtain a specialized PIOP whose prover complexity $O(NMn)$ is linear in the witness size and a $\log n$ factor better than any approach that treats ring multiplication as a black-box circuit gate.  A further practical advantage of working over splitting fields (rather than Galois rings) is compatibility with the DP25 packing reduction~\cite{EC:DiaPos25}: since our sumcheck and oracles live over the extension field $\extensionfield$, which admits an explicit tower $\field_p \subset \field_{p^2} \subset \cdots \subset \extensionfield$, we can apply the $s$-round dimension reduction to compress PCS costs. The Galois-ring systems cannot directly exploit this tower structure, as their native domain is a ring rather than a field.

%%%%%%%%%%%%%%%%%%%%%%%%%%%%%%%%%%%%%%%%%%%%%%%%%%%%%%%%%%%%%%%%%%%%%%%%%%%%%%%
\paragraph{Ring-switching proofs over polynomial rings.}
Huang, Mao, and Zhang~\cite{EPRINT:HuaMaoZha25} propose a proof system for Ring-R1CS over prime-power modulus cyclotomic rings.
%They address two obstacles: the exceptional set of $\ZZ_q$ may be too small to achieve soundness without a large number of repetitions, and when the proof elements are over (extensions of) the ring they yield a multiplicative factor of $n$ in the proof size.
Their key technique is \emph{ring switching}: for each row $i$ of $As = t$, the prover computes the quotient $Q_i(X) = \bigl(\sum_j A_{ij}(X) \cdot s_j(X) - t_i(X)\bigr)/(X^n+1)$, lifting the mod-$(X^n+1)$ congruence to an exact identity in $\ZZ_q[X]$. Assuming preprocessed $A$, they would still send oracles for $\tilde{s} \in (\field_{q}^{<n}[X])[Y_1,\dotsc,Y_{\log M}]$ and $\tilde{Q}\in (\field_{q}^{<3n}[X])[Y_1,\dotsc,Y_{\log N}]$. They consider the extension field $\extensionfield = \ZZ_{q^k}$ such that $q^k \geq 2^\lambda$. The identity is evaluated at a random $X = \alpha \in K$, collapsing each ring equation to a scalar identity in $\extensionfield$, and then a standard multilinear sumcheck is used to reduce this to verifying PCS openings.

While in some ways our protocols are similar, we believe our scheme presents a few benefits. First, since our scheme works with the cyclotomic structure in the sumcheck argument machinery rather than collapsing it, we don't need to compute or send $\tilde{Q}$, reducing commitment costs from $O(nNM)$ to $O(nM)$ in the holographic setting. In their full quadratic R1CS protocol $Q$'s degree in $X$ becomes $3N$; practically speaking, since $M \approxeq N$ for R1CS circuits their protocol requires committing to $4\times$ the base field elements\footnote{Empirically, $M/N \approx 1$ across diverse practical R1CS instantiations, including both algebraically-oriented circuits (Aptos keyless OpenID, \url{https://github.com/aptos-labs/keyless-zk-proofs}) and bitwise-dominated circuits (circomlib SHA-256, \url{https://github.com/iden3/circomlib/blob/master/circuits/sha256/sha256.circom}).}. Second, typing oracles over Galois rings $\mathsf{GR}_{q}^{n}(\field_{q}^{<n}[X])$ is non-standard and has required extending many theorems about codes~\cite{ESORICS:WZDWZY25,C:JLXYY25}; these results, and thus the efficiency of the resulting polynomial commitments, still lag behind those for finite fields. Hachi~\cite{hachi} introduces a generic reduction from polynomial evaluation over extension fields to cyclotomic rings (so this is relevant to our PIOP as well), but state-of-the-art lattice PCs still lag behind their hash-based counterparts. At the information-theoretic level, it is unknown whether many PIORs can be applied. One important example of this is the ring-switch based \emph{packing} technique of Diamond and Posen~\cite{EC:DiaPos25} that we show can be applied in our setting to reduce $\hat{s}$ to just $\log M - \log [\extensionfield : \field_p]$ variables.

%%%%%%%%%%%%%%%%%%%%%%%%%%%%%%%%%%%%%%%%%%%%%%%%%%%%%%%%%%%%%%%%%%%%%%%%%%%%%%%
\paragraph{Works proving FHE.}

Laminate~\cite{EPRINT:PLFT25} presents a verifiable FHE construction. It proves a layered payload circuit via a GKR-style blind hIOP while preserving BFV/BGV's native SIMD layout; the protocol can be understood as amortizing over running a proof for each CRT slot in parallel, using standard extension field techniques to amplify soundness, whereas our commitments are in the coefficient basis.

Wei et al.~\cite{EPRINT:WWXZDZ25} give a new sumcheck argument over fields of small characteristic that asymptotically reduces prover multiplications, while requiring sending $2kd$ oracles in each of $dk$ rounds. As such, it remains impractical.

%%%%%%%%%%%%%%%%%%%%%%%%%%%%%%%%%%%%%%%%%%%%%%%%%%%%%%%%%%%%%%%%%%%%%%%%%%%%%%%
\paragraph{DP25 ring-switch packing.}
Diamond and Posen~\cite{EC:DiaPos25} introduce a packing technique that exploits the tower of quadratic extensions $\field_p \subset \field_{p^2} \subset \cdots \subset \extensionfield$ to compress multilinear polynomial dimensions by $s = \log [\extensionfield : \field_p]$ via an $s$-round degree-1 sumcheck. We use this as a post-processing step (\cref{sec:reducing-commitments}) to reduce the polynomial commitment cost, and we provide an odd-prime tower construction (\cref{sec:odd-prime-tower}) that extends DP25's original even-characteristic instantiation.

As mentioned before, DP25 packing requires a tower of \emph{field} extensions with an explicit multilinear basis at each level. The Galois-ring proof systems~\cite{C:JLXYY25,PKC:WeiZhaDen25,ESORICS:WZDWZY25} work over the ring $\mathrm{GR}(p^k, d) = \ZZ_{p^k}[X]/(h(X))$, so DP25's tower construction does not directly apply. These systems instead use the canonical embedding $\ZZ_{p^k} \hookrightarrow \mathrm{GR}(p^k, d)$.